% B05 / A023 complete positive-credit bound for disjoint point defects
\section*{A uniform single-use payment for finite point-defect products}

Let $q_n$ be $+1$ at one point of the uniform $n$-cube and $-1$
elsewhere.  Consider a finite product of such factors on pairwise disjoint
physical blocks.  For a collection of unrevealed factors with disagreement
probability $y\in[0,1/2]$, the complete conditional-label plus phase packet
against a constant is $F(y)$, where
\[
 F(p)=4p(1-p)-2(1-p)^2\log_2(1-p)-2p^2\log_2p
 +(1-2p)^2\log_2(1-2p),\qquad 0\le p\le\tfrac12.       \tag{1}
\]
The endpoint convention is $0\log 0=0$.

Fix any deterministic factor order and exhaust the current factor $q_a$
before moving to the next one.  Put $\delta=2^{-a}$.  At residual dimension
$m$, with $x=2^{-m}$, define
\[
 \varphi_y(x)=F\bigl(y+(1-2y)x\bigr).
\]
For $m\ge2$ the reached local credit is
\[
 \Gamma_m(y)=\varphi_y(x)-\frac{\varphi_y(2x)+\varphi_y(0)}2.           \tag{2}
\]
Writing
\[
 s_y(x)=\frac{\varphi_y(x)-\varphi_y(0)}x,
\]
equation (2) becomes the exact dyadic secant identity
\[
 \Gamma_m(y)=x\bigl(s_y(x)-s_y(2x)\bigr).                              \tag{3}
\]
The node has reach probability $2^{m-a}=\delta/x$, so its reached positive
credit is
\[
 \delta\,[s_y(x)-s_y(2x)]_+.                                          \tag{4}
\]

The needed shape control is global and elementary.  Differentiation of (1)
gives, on $(0,1/2)$,
\[
 F''(p)=-8+\frac4{\ln2}\ln\frac{(1-2p)^2}{p(1-p)},
 \qquad
 F'''(p)=-\frac4{(\ln2)p(1-p)(1-2p)}<0.                               \tag{5}
\]
Thus $\varphi_y''$ is strictly decreasing unless $y=1/2$, in which case
$\varphi_y$ is constant.  Moreover
\[
 s_y'(x)=\frac1{x^2}\int_0^x t\varphi_y''(t)\,dt.                     \tag{6}
\]
The integral in (6) first increases and then decreases, and crosses zero at
most once.  Hence $s_y$ is unimodal: it is first nondecreasing and then
nonincreasing.  This includes $y=0$ by the integrable endpoint limit.

Also
\[
 F'(p)=4-8p+\frac4{\ln2}igl[(1-p)\ln(1-p)-p\ln p
 -(1-2p)\ln(1-2p)\bigr].                                              \tag{7}
\]
Since $|t\ln t|\le1$ on $[0,1]$ and
$\ln2=2(1/3+1/(3^3\cdot3)+\cdots)>2/3$, equation (7) gives
\[
 |F'(p)|<22.                                                           \tag{8}
\]
Therefore $|s_y(x)|<22$.  Along the decreasing dyadic sequence
$1/2,1/4,\ldots,2^{-a}$, unimodality implies that the sum of all positive
increments in (4) is less than the full range $44$.  Consequently the
contribution of all levels $m\ge2$ is less than $44\delta$.

At $m=1$, both children have packet $F(y)$, while the parent has packet
$F(1/2)=2$.  Nonnegativity of the complete packet gives the separate bound
\[
 2^{1-a}[F(1/2)-F(y)]_+\le4\delta.                                    \tag{9}
\]
Combining (4)--(9), the total reached positive credit $P(a\mid y)$ while
exhausting this factor satisfies
\[
 \boxed{P(a\mid y)<48\,2^{-a}.}                                      \tag{10}
\]

The total influence of $q_a$ is
\[
 I(q_a)=2a,2^{-a}\ge2\,2^{-a}.
\]
Thus (10) implies
\[
 P(a\mid y)<24I(q_a).                                                  \tag{11}
\]
Summing (11) in any factor order and using exact disjoint-product influence
additivity proves
\[
 \boxed{\sum_jP_j<24\sum_jI(q_{n_j})=24I\!\left(\prod_jq_{n_j}\right).}\tag{12}
\]

The preceding ordered M\"obius theorem gives the exact single-use identity
\[
 R(A)+\sum_jN_j=\sum_jP_j.                                             \tag{13}
\]
Equations (12)--(13) are therefore a complete finite physical stopping-tree
payment for every finite disjoint product of point defects:
\[
 \boxed{R(A)+\sum_jN_j<24I\!\left(\prod_jq_{n_j}\right).}              \tag{14}
\]
Every local credit is used once, and there is no terminal packet because all
factors are exhausted.  The order is arbitrary, not optimized.

This is a restricted family theorem.  It does not cover unrelated addressing
branches, shared physical variables, or arbitrary Boolean leaves, and it does
not prove the full Fourier Entropy--Influence conjecture.
